% songs/es-war-an-einem-sommertag.tex
\begin{song}{Es war an einem Sommertag}{Arno Claus}{Capo -}

\slbl{Intro:}
\cl{C  G  C  G  Am  E  Am}

\cl{   C      G      C     G    Am            E     Am}
\slbl{1.}Es war an einem Sommertag, irgendwann und irgendwo,
\cl{   C      G        C          G       Am     E         Am}
Da tönte plötzlich Trommelschlag und Flötenspiel klang froh.
\cl{    G               C                G                      C}
Es war eine muntre bunte Schar, die machte vor dem Rathaus Halt,
\cl{     G                  C            Am           E     Am}
Den Grund, weshalb sie gekommen war, erfuhren die Leute bald.

\slbl{Interlude:}
\cl{C  G  C  G  Am  E  Am}

\cl{    C        G     C     G        Am       E     Am}
\slbl{2.}Ein Mann mit einem Federhut rief: Männer, hört mir zu!
\cl{C      G            C        G       Am    E    Am}
Ich versprech' euch Geld und Gut, und Ehre noch dazu!
\cl{    G                    C                  G                      C}
Der Kaiser braucht euch, reiht euch ein, hängt nicht an Weib und Haus!
\cl{     G                  C                Am       E     Am}
Es wird auch gar nicht lange sein, zieht mit ins Feld hinaus!

\slbl{Interlude:}
\cl{C  G  C  G  Am  E  Am}

\cl{    C        G        C      G        Am           E       Am}
\slbl{3.}Im Wirtshaus war das Trinken frei, bezahlt mit des Kaisers Gold.
\cl{     C       G      C    G       Am              E      Am}
Und während dieser Zecherei trat mancher in des Kaisers Sold.
\cl{    G               C              G                        C}
Gab seiner Frau den Abschiedskuß, versuchte als Soldat sein Glück,
\cl{     G                  C            Am        E     Am}
sah nicht des Werbers Pferdefuß und kehrt' nie mehr zurück.

\slbl{Interlude:}
\cl{C  G  C  G  Am  E  Am}

\cl{    C        G        C      G           Am      E       Am}
\slbl{4.}Mit Flötenspiel und Trommelschlag ging's früh am Morgen fort.
\cl{    C           G        C      G     Am          E       Am}
Die Schar ward größer, denn es lag am Weg noch so mancher Ort.
\cl{    G               C              G                        C}
Der Werber mit dem Federhut macht' sein Geschäft nicht schlecht,
\cl{     G                  C               Am       E        Am}
versprach noch vielen Geld und Gut, dem Kaiser, dem war's recht.

\slbl{Interlude:}
\cl{C  G  C  G  Am  E  Am}

\cl{    C      G     C      G         Am      E       Am}
\slbl{5.}Die Jahre gingen in das Land, und von der großen Schar,
\cl{      C       G        C      G        Am    E      Am}
gab's keinen, der nach Hause fand, wie er gegeangen war.
\cl{    G               C            G                     C}
Der eine ließ sein Bein im Feld, blind karm ein andrer an.
\cl{     G                  C               Am     E      Am}
Die meisten hatte der Tod gefällt, der jede Schlacht gewann.

\slbl{Interlude:}
\cl{C  G  C  G  Am  E  Am}

\cl{     C      G      C     G     Am         E     Am}
\slbl{6.}Die letzten Tränen waren kaum geweint, da waren sie,
\cl{      C        G      C        G         Am        E     Am}
auch schon vergessen, wie ein Traum, die Menschen lernen nie.
\cl{    G               C           G                  C}
Und dann, an einem Sommertag, irgendwann und irgendwo,
\cl{     G                  C              Am     E         Am}
da tönte plötzlich Trommelschlag, und Flötenspiel klang froh.

\slbl{Outro:}
\cl{C  G  C  G  Am  E  Am}

\end{song}